\section{Background}
\label{sec:background}

\subsection{The MiniMind Architecture}
\label{sec:background:arch}

MiniMind implements a Llama-family~\citep{touvron2023llama} transformer decoder
with the following design choices, matching the Qwen3~\citep{qwen3} model family:

\begin{itemize}[leftmargin=*]
  \item \textbf{RMSNorm} pre-normalization with $\epsilon = 10^{-6}$.
  \item \textbf{Grouped-query attention (GQA)}~\citep{ainslie2023gqa}: 8 query
    heads, 4 key-value heads, per-head RMSNorm on $Q$ and $K$.
  \item \textbf{Rotary position embeddings (RoPE)}~\citep{su2024roformer} with
    $\theta = 10^6$.  Inference-time YaRN~\citep{peng2023yarn} scaling (4$\times$
    extension) is available via \texttt{--inference\_rope\_scaling}.
  \item \textbf{SwiGLU} feed-forward network; intermediate size derived as
    $\lceil h \cdot \pi / 64 \rceil \times 64$ where $h$ is \texttt{hidden\_size}.
  \item \textbf{Flash attention}: \texttt{F.scaled\_dot\_product\_attention}
    (PyTorch SDPA), which dispatches to cuDNN-FlashAttention on CUDA and to
    hipBLASLt-fused kernels on ROCm.  Falls back to manual attention when
    \texttt{seq\_len=1} (decode step) or when a padding mask is present.
  \item \textbf{Tied embeddings} (\texttt{tie\_word\_embeddings=True});
    vocabulary size 6400, trained from a custom BPE tokenizer.
\end{itemize}

\paragraph{Dense configuration.}
Default: \texttt{hidden\_size=768}, \texttt{num\_hidden\_layers=8},
\texttt{num\_attention\_heads=8}, \texttt{num\_key\_value\_heads=4}.
Total parameters: $\approx$64\,M.

\paragraph{MoE configuration.}
A sparse mixture-of-experts~\citep{shazeer2017outrageously} variant replaces
the dense FFN with a gated MoE layer: 4 routed experts, 1 active per token, plus
a shared expert, governed by
\texttt{router\_aux\_loss\_coef}$= 5 \times 10^{-4}$.  Total parameters:
$\approx$198\,M; active per forward: $\approx$64\,M.  The loss function always
computes \texttt{loss = res.loss + res.aux\_loss} so the router
load-balancing term is included during both dense and MoE training
(it is identically zero for dense models).

\paragraph{Tokenizer.}
A BPE tokenizer with vocabulary size 6400, trained on the project's text corpus.
Special tokens include \texttt{<think>}, \texttt{</think>},
\texttt{<tool\_call>}, \texttt{<tool\_response>}, enabling chain-of-thought and
tool-use workflows directly in the chat template.

\subsection{Algorithm Landscape}
\label{sec:background:algos}

MiniMind covers the following training stages, each implemented as a self-contained
script in \texttt{trainer/}:

\paragraph{Pretraining.}
Standard causal language modeling loss over next-token prediction.  The trainer
(\texttt{train\_pretrain.py}) supports DDP via \texttt{torchrun}, bf16/fp16 mixed
precision via \texttt{torch.amp}, cosine learning-rate schedule, gradient
clipping, and optional \texttt{torch.compile}.

\paragraph{Supervised Fine-Tuning (SFT).}
\texttt{train\_full\_sft.py} uses \texttt{tokenizer.apply\_chat\_template} to
format instruction-response pairs and masks all non-assistant tokens
(labels $= -100$), so the loss is computed only on assistant completions.

\paragraph{LoRA.}
\texttt{train\_lora.py} injects low-rank adapters via a custom
\texttt{apply\_lora} function in \texttt{model/model\_lora.py} (no \texttt{peft}
dependency).  Only LoRA parameters are trained; all base-model weights are
frozen.  The adapter is saved as a separate \texttt{.pth} keyed by layer name.
\texttt{torch.compile} is automatically disabled for LoRA (monkey-patch forward
incompatibility is caught and logged).

\paragraph{Knowledge Distillation.}
\texttt{train\_distillation.py} minimizes a weighted sum of cross-entropy (CE)
loss against ground-truth labels and Kullback-Leibler (KL) divergence against
a frozen teacher's logits:
\[
  \mathcal{L} = \alpha \mathcal{L}_{\text{CE}} + (1 - \alpha)
    \cdot T^2 \cdot D_{\text{KL}}(p_{\text{teacher}} \,\|\, p_{\text{student}})
\]
where $T$ is the temperature and $\alpha \in [0,1]$ controls the CE weight.
The default configuration distills from a MoE teacher into a dense student.

\paragraph{Direct Preference Optimization (DPO).}
\texttt{train\_dpo.py} implements DPO~\citep{rafailov2023direct} with a frozen
reference model.  The loss is:
\[
  \mathcal{L}_{\text{DPO}} = -\mathbb{E}_{(x, y_w, y_l)} \left[
    \log \sigma \!\left( \beta \left(
      \log \frac{\pi_\theta(y_w|x)}{\pi_{\text{ref}}(y_w|x)} -
      \log \frac{\pi_\theta(y_l|x)}{\pi_{\text{ref}}(y_l|x)}
    \right) \right)
  \right]
\]
with $\beta = 0.15$ (default) and a very small learning rate ($4 \times 10^{-8}$)
to prevent catastrophic forgetting.

\paragraph{Proximal Policy Optimization (PPO).}
\texttt{train\_ppo.py} maintains four models simultaneously: an actor, a critic
(value model), a frozen reference, and a separately loaded reward model
(InternLM2-1.8B-reward by default).  Rollouts are collected by
\texttt{rollout\_engine.py} (\texttt{--rollout\_engine torch} or
\texttt{--rollout\_engine sglang}).  The actor is updated using clipped
PPO~\citep{schulman2017proximal} with generalized advantage estimation (GAE),
early-stopped when per-step KL exceeds a threshold.

\paragraph{Group-Relative Policy Optimization (GRPO).}
\texttt{train\_grpo.py} implements GRPO~\citep{shao2024deepseekmath}: for each
prompt, $G$ responses are sampled; rewards are computed by a reward model plus
heuristic rule terms (response length, \texttt{<think>} format, repetition
penalty); advantages are normalized within the group.  No critic model is needed.
The thinking ratio ($\texttt{--thinking\_ratio}$) controls how often the
\texttt{open\_thinking} chat-template flag is activated.

\paragraph{Agent RL.}
\texttt{train\_agent.py} extends GRPO to multi-turn agent trajectories involving
tool calls; the reward function evaluates both the quality of tool invocations and
the final answer.

\subsection{Hardware and Runtime Stacks}
\label{sec:background:hardware}

\begin{table}[h]
\centering
\caption{Hardware and software stacks under test.}
\label{tab:hw}
\small
\begin{tabular}{lll}
\toprule
\textbf{Axis} & \textbf{NVIDIA} & \textbf{AMD} \\
\midrule
GPU & H100 SXM5 (80\,GB HBM3) & MI350 (HBM3e) \\
\# GPUs / node & 8 & 8 \\
Intra-node interconnect & NVLink 4 + NVSwitch & Infinity Fabric \\
PyTorch wheel & \texttt{2.10.0+cu12x} & \texttt{2.10.0+rocm7.1} \\
Distributed backend & NCCL & RCCL (aliased as \texttt{"nccl"}) \\
Compiler stack & CUDA / cuDNN / cuBLAS / Triton & ROCm / MIOpen / hipBLASLt / Triton \\
Low-precision & bf16, fp16; Transformer Engine FP8 & bf16, fp16; ROCm FP8 \\
Inference engines & HF native, vLLM (CUDA), SGLang (CUDA) & HF native, vLLM-ROCm, SGLang-ROCm \\
\bottomrule
\end{tabular}
\end{table}

\paragraph{PyTorch portability.}
PyTorch-ROCm aliases the \texttt{torch.cuda.*} API surface to HIP so that code
referencing \texttt{device\_type="cuda"}, \texttt{torch.cuda.is\_available()},
or \texttt{torch.cuda.get\_device\_properties()} works unchanged on AMD GPUs.
MiniMind's \texttt{device\_utils.py} detects the ROCm build by inspecting
\texttt{torch.version.hip} and adjusts logging accordingly; the rest of the
codebase is unaware of the vendor.

\paragraph{SDPA dispatch.}
\texttt{F.scaled\_dot\_product\_attention} selects its backend at runtime.  On
CUDA, it may dispatch to cuDNN FlashAttention, FlashAttention-2, or
FlashAttention-3 depending on the cuDNN version and \texttt{sdp\_kernel} context.
On ROCm, it dispatches to hipBLASLt-fused attention or Composable Kernel (CK)
flash attention.  We report which backend was selected for each datapoint via
\texttt{torch.backends.cuda.flash\_sdp\_enabled()} and its ROCm equivalent.

\paragraph{Collective communication.}
Under DDP, \texttt{dist.init\_process\_group("nccl")} maps to NCCL on NVIDIA and
RCCL on AMD.  Both expose the same \texttt{torch.distributed} API.  Intra-node
bandwidth (all-reduce of the $\approx$64\,M-parameter gradient buffer at bf16)
is reported in~\S\ref{sec:crosscutting:nccl}.
